\Chapter{118}

\Verse{1}
{
    \zR{话说邢王二夫人听尤氏一段话，明知也难挽回。}
}
{
    \eR{[English source not present in the checked-in Bencraft/Joly files.]}
}
{
}

\Verse{2}
{
    \zR{王夫人只得说道：“姑娘要行 善，这也是前生的夙根，我们也实在拦不住。}
}
{
    \eR{[English source not present in the checked-in Bencraft/Joly files.]}
}
{
}

\Verse{3}
{
    \zR{只是咱们这样人家的姑娘出了家，不 成个事体。}
}
{
    \eR{[English source not present in the checked-in Bencraft/Joly files.]}
}
{
}

\Verse{4}
{
    \zR{如今你嫂子说了，准你修行，也是好处。}
}
{
    \eR{[English source not present in the checked-in Bencraft/Joly files.]}
}
{
}

\Verse{5}
{
    \zR{却有一句话要说：那头发可以 不剃的，只要自己的心真，那在头发上头呢？}
}
{
    \eR{[English source not present in the checked-in Bencraft/Joly files.]}
}
{
}

\Verse{6}
{
    \zR{你想妙玉也是带发修行的。}
}
{
    \eR{[English source not present in the checked-in Bencraft/Joly files.]}
}
{
}

\Verse{7}
{
    \zR{不知他怎 样凡心一动，才闹到那个分儿。}
}
{
    \eR{[English source not present in the checked-in Bencraft/Joly files.]}
}
{
}

\Verse{8}
{
    \zR{姑娘执意如此，我们就把姑娘住的房子便算了姑娘 的静室。}
}
{
    \eR{[English source not present in the checked-in Bencraft/Joly files.]}
}
{
}

\Verse{9}
{
    \zR{所有服侍姑娘的人，也得叫他们来问。}
}
{
    \eR{[English source not present in the checked-in Bencraft/Joly files.]}
}
{
}

\Verse{10}
{
    \zR{他若愿意跟的，就讲不得说亲配人； 若不愿意跟的，另打主意。”}
}
{
    \eR{[English source not present in the checked-in Bencraft/Joly files.]}
}
{
}

\Verse{11}
{
    \zR{惜春听了，收了泪，拜谢了邢王二夫人，李纨、尤氏 等。}
}
{
    \eR{[English source not present in the checked-in Bencraft/Joly files.]}
}
{
}

\Verse{12}
{
    \zR{王夫人说了，便问彩屏等：“谁愿跟姑娘修行？”}
}
{
    \eR{[English source not present in the checked-in Bencraft/Joly files.]}
}
{
}

\Verse{13}
{
    \zR{彩屏等回道：“太太们派谁就 是谁。”}
}
{
    \eR{[English source not present in the checked-in Bencraft/Joly files.]}
}
{
}

\Verse{14}
{
    \zR{王夫人知道不愿意，正在想人。}
}
{
    \eR{[English source not present in the checked-in Bencraft/Joly files.]}
}
{
}

\Verse{15}
{
    \zR{袭人立在宝玉身后，想来宝玉必要大哭，防着 他的旧病。}
}
{
    \eR{[English source not present in the checked-in Bencraft/Joly files.]}
}
{
}

\Verse{16}
{
    \zR{岂知宝玉叹道：“真真难得！”}
}
{
    \eR{[English source not present in the checked-in Bencraft/Joly files.]}
}
{
}

\Verse{17}
{
    \zR{袭人心里更自伤悲。}
}
{
    \eR{[English source not present in the checked-in Bencraft/Joly files.]}
}
{
}

\Verse{18}
{
    \zR{宝钗虽不言语，遇事 试探，见他执迷不醒，只得暗中落泪。}
}
{
    \eR{[English source not present in the checked-in Bencraft/Joly files.]}
}
{
}

\Verse{19}
{
    \zR{王夫人才要叫了众丫头来问，忽见紫鹃走上 前去，在王夫人面前跪下，回道：“刚才太太问跟四姑娘的姐姐，太太看着怎么样？”}
}
{
    \eR{[English source not present in the checked-in Bencraft/Joly files.]}
}
{
}

\Verse{20}
{
    \zR{王夫人道：“这个如何强派得人的？}
}
{
    \eR{[English source not present in the checked-in Bencraft/Joly files.]}
}
{
}

\Verse{21}
{
    \zR{谁愿意，他自然就说出来了。”}
}
{
    \eR{[English source not present in the checked-in Bencraft/Joly files.]}
}
{
}

\Verse{22}
{
    \zR{紫鹃道：“姑娘修 行，自然姑娘愿意，并不是别的姐姐们的意思。}
}
{
    \eR{[English source not present in the checked-in Bencraft/Joly files.]}
}
{
}

\Verse{23}
{
    \zR{我有句话回太太：我也并不是拆开 姐姐们，各人有各人的心。}
}
{
    \eR{[English source not present in the checked-in Bencraft/Joly files.]}
}
{
}

\Verse{24}
{
    \zR{我服侍林姑娘一场，林姑娘待我也是太太们知道的，实 在恩重如山，无以可报。}
}
{
    \eR{[English source not present in the checked-in Bencraft/Joly files.]}
}
{
}

\Verse{25}
{
    \zR{他死了，我恨不得跟了他去，但只他不是这里的人，我又 受主子家的恩典，难以从死。}
}
{
    \eR{[English source not present in the checked-in Bencraft/Joly files.]}
}
{
}

\Verse{26}
{
    \zR{如今四姑娘既要修行，我就求太太们将我派了跟着姑 娘，伏侍姑娘一辈子，不知太太们准不准？}
}
{
    \eR{[English source not present in the checked-in Bencraft/Joly files.]}
}
{
}

\Verse{27}
{
    \zR{若准了，就是我的造化了。”}
}
{
    \eR{[English source not present in the checked-in Bencraft/Joly files.]}
}
{
}

\Verse{28}
{
    \zR{邢王二夫人 尚未答言，只见宝玉听到那里，想起黛玉，一阵心酸，眼泪早下来了。}
}
{
    \eR{[English source not present in the checked-in Bencraft/Joly files.]}
}
{
}

\Verse{29}
{
    \zR{众人才要问 他时，他又哈哈的大笑，走上来道：“我不该说的。}
}
{
    \eR{[English source not present in the checked-in Bencraft/Joly files.]}
}
{
}

\Verse{30}
{
    \zR{这紫鹃蒙太太派给我屋里，我 才敢说：求太太准了他罢，全了他的好心。”}
}
{
    \eR{[English source not present in the checked-in Bencraft/Joly files.]}
}
{
}

\Verse{31}
{
    \zR{王夫人道：“你头里姊妹出了嫁，还哭 得死去活来；如今看见四妹妹要出家，不但不劝，倒说‘好事’。}
}
{
    \eR{[English source not present in the checked-in Bencraft/Joly files.]}
}
{
}

\Verse{32}
{
    \zR{你如今到底是怎 么个意思？}
}
{
    \eR{[English source not present in the checked-in Bencraft/Joly files.]}
}
{
}

\Verse{33}
{
    \zR{我索性不明白了。”}
}
{
    \eR{[English source not present in the checked-in Bencraft/Joly files.]}
}
{
}

\Verse{34}
{
    \zR{宝玉道：“四妹妹修行是已经准了的，四妹妹也是一 定的主意了？}
}
{
    \eR{[English source not present in the checked-in Bencraft/Joly files.]}
}
{
}

\Verse{35}
{
    \zR{若是真呢，我有一句话告诉太太；若是不定呢，我就不敢混说了。”}
}
{
    \eR{[English source not present in the checked-in Bencraft/Joly files.]}
}
{
}

\Verse{36}
{
    \zR{惜 春道：“二哥哥说话也好笑，一个人主意不定，便扭得过太太们来了？}
}
{
    \eR{[English source not present in the checked-in Bencraft/Joly files.]}
}
{
}

\Verse{37}
{
    \zR{我也是像紫鹃 的话：容我呢，是我的造化；不容我呢还有一个死呢：那怕什么？}
}
{
    \eR{[English source not present in the checked-in Bencraft/Joly files.]}
}
{
}

\Verse{38}
{
    \zR{二哥哥既有话， 只管说。”}
}
{
    \eR{[English source not present in the checked-in Bencraft/Joly files.]}
}
{
}

\Verse{39}
{
    \zR{宝玉道：“我这也不算什么泄漏了，这也是一定的。}
}
{
    \eR{[English source not present in the checked-in Bencraft/Joly files.]}
}
{
}

\Verse{40}
{
    \zR{我念一首诗给你们听 听罢。”}
}
{
    \eR{[English source not present in the checked-in Bencraft/Joly files.]}
}
{
}

\Verse{41}
{
    \zR{众人道：“人家苦得很的时候，你倒来做诗怄人。”}
}
{
    \eR{[English source not present in the checked-in Bencraft/Joly files.]}
}
{
}

\Verse{42}
{
    \zR{宝玉道：“不是做诗，我 到过一个地方儿看了来的。}
}
{
    \eR{[English source not present in the checked-in Bencraft/Joly files.]}
}
{
}

\Verse{43}
{
    \zR{你们听听罢。”}
}
{
    \eR{[English source not present in the checked-in Bencraft/Joly files.]}
}
{
}

\Verse{44}
{
    \zR{众人道：“使得。}
}
{
    \eR{[English source not present in the checked-in Bencraft/Joly files.]}
}
{
}

\Verse{45}
{
    \zR{你就念念，别顺着嘴儿 胡诌。”}
}
{
    \eR{[English source not present in the checked-in Bencraft/Joly files.]}
}
{
}

\Verse{46}
{
    \zR{宝玉也不分辩，便说道： 勘破三春景不长，缁衣顿改昔年妆。}
}
{
    \eR{[English source not present in the checked-in Bencraft/Joly files.]}
}
{
}

\Verse{47}
{
    \zR{可怜绣户侯门女，独卧青灯古佛旁。}
}
{
    \eR{[English source not present in the checked-in Bencraft/Joly files.]}
}
{
}

\Verse{48}
{
    \zR{李纨宝钗听了，诧异道：“不好了!这个人入了魔了。”}
}
{
    \eR{[English source not present in the checked-in Bencraft/Joly files.]}
}
{
}

\Verse{49}
{
    \zR{王夫人听了这话，点头叹息， 便问：“宝玉，你到底是那里看来的？”}
}
{
    \eR{[English source not present in the checked-in Bencraft/Joly files.]}
}
{
}

\Verse{50}
{
    \zR{宝玉不便说出来，回道：“太太也不必问我， 自有见的地方。”}
}
{
    \eR{[English source not present in the checked-in Bencraft/Joly files.]}
}
{
}

\Verse{51}
{
    \zR{王夫人回过味来，细细一想，便更哭起来道：“你说前儿是玩话， 怎么忽然有这首诗？}
}
{
    \eR{[English source not present in the checked-in Bencraft/Joly files.]}
}
{
}

\Verse{52}
{
    \zR{罢了，我知道了。}
}
{
    \eR{[English source not present in the checked-in Bencraft/Joly files.]}
}
{
}

\Verse{53}
{
    \zR{你们叫我怎么样呢？}
}
{
    \eR{[English source not present in the checked-in Bencraft/Joly files.]}
}
{
}

\Verse{54}
{
    \zR{我也没有法儿了，也只得 由着你们去罢，但只等我合上了眼，各自干各自的就完了！”}
}
{
    \eR{[English source not present in the checked-in Bencraft/Joly files.]}
}
{
}

\Verse{55}
{
    \zR{宝钗一面劝着，这个心比刀绞更甚，也掌不住，便放声大哭起来。}
}
{
    \eR{[English source not present in the checked-in Bencraft/Joly files.]}
}
{
}

\Verse{56}
{
    \zR{袭人已经哭 的死去活来，幸亏秋纹扶着。}
}
{
    \eR{[English source not present in the checked-in Bencraft/Joly files.]}
}
{
}

\Verse{57}
{
    \zR{宝玉也不啼哭，也不相劝，只不言语。}
}
{
    \eR{[English source not present in the checked-in Bencraft/Joly files.]}
}
{
}

\Verse{58}
{
    \zR{贾兰贾环听到 那里，各自走开。}
}
{
    \eR{[English source not present in the checked-in Bencraft/Joly files.]}
}
{
}

\Verse{59}
{
    \zR{李纨竭力的解说：“总是宝兄弟见四妹妹修行，他想来是痛极了， 不顾前后的疯话，这也作不得准。}
}
{
    \eR{[English source not present in the checked-in Bencraft/Joly files.]}
}
{
}

\Verse{60}
{
    \zR{独有紫鹃的事情，准不准，好叫他起来。”}
}
{
    \eR{[English source not present in the checked-in Bencraft/Joly files.]}
}
{
}

\Verse{61}
{
    \zR{王夫 人道：“什么依不依？}
}
{
    \eR{[English source not present in the checked-in Bencraft/Joly files.]}
}
{
}

\Verse{62}
{
    \zR{横竖一个人的主意定了，那也是扭不过来的。}
}
{
    \eR{[English source not present in the checked-in Bencraft/Joly files.]}
}
{
}

\Verse{63}
{
    \zR{可是宝玉说的， 也是一定的了！”}
}
{
    \eR{[English source not present in the checked-in Bencraft/Joly files.]}
}
{
}

\Verse{64}
{
    \zR{紫鹃听了磕头。}
}
{
    \eR{[English source not present in the checked-in Bencraft/Joly files.]}
}
{
}

\Verse{65}
{
    \zR{惜春又谢了王夫人。}
}
{
    \eR{[English source not present in the checked-in Bencraft/Joly files.]}
}
{
}

\Verse{66}
{
    \zR{紫鹃又给宝玉宝钗磕了头， 宝玉念声：“阿弥陀佛!难得，难得!不料你倒先好了。”}
}
{
    \eR{[English source not present in the checked-in Bencraft/Joly files.]}
}
{
}

\Verse{67}
{
    \zR{宝钗虽然有把持，也难掌住。}
}
{
    \eR{[English source not present in the checked-in Bencraft/Joly files.]}
}
{
}

\Verse{68}
{
    \zR{只有袭人也顾不得王夫人在上，便痛哭不止，说：“我也愿意跟了四姑娘去修行。”}
}
{
    \eR{[English source not present in the checked-in Bencraft/Joly files.]}
}
{
}

\Verse{69}
{
    \zR{宝玉笑道：“你也是好心，但是你不能享这个清福的。”}
}
{
    \eR{[English source not present in the checked-in Bencraft/Joly files.]}
}
{
}

\Verse{70}
{
    \zR{袭人哭道：“这么说，我是 要死的了？”}
}
{
    \eR{[English source not present in the checked-in Bencraft/Joly files.]}
}
{
}

\Verse{71}
{
    \zR{宝玉听到那里，倒觉伤心，只是说不出来。}
}
{
    \eR{[English source not present in the checked-in Bencraft/Joly files.]}
}
{
}

\Verse{72}
{
    \zR{因时已五更，宝玉请王夫人安歇。}
}
{
    \eR{[English source not present in the checked-in Bencraft/Joly files.]}
}
{
}

\Verse{73}
{
    \zR{李纨等各自散去。}
}
{
    \eR{[English source not present in the checked-in Bencraft/Joly files.]}
}
{
}

\Verse{74}
{
    \zR{彩屏等暂且伏侍惜春回去， 后来指配了人家。}
}
{
    \eR{[English source not present in the checked-in Bencraft/Joly files.]}
}
{
}

\Verse{75}
{
    \zR{紫鹃终身伏侍，毫不改初。}
}
{
    \eR{[English source not present in the checked-in Bencraft/Joly files.]}
}
{
}

\Verse{76}
{
    \zR{此是后话。}
}
{
    \eR{[English source not present in the checked-in Bencraft/Joly files.]}
}
{
}

\Verse{77}
{
    \zR{且言贾政扶了贾母灵柩，一路南行，因遇着班师的兵将船只过境，河道拥挤， 不能速行，在道实在心焦。}
}
{
    \eR{[English source not present in the checked-in Bencraft/Joly files.]}
}
{
}

\Verse{78}
{
    \zR{幸喜遇见了海疆的官员，闻得镇海统制钦召回京，想来 探春一定回家，略略解些烦心。}
}
{
    \eR{[English source not present in the checked-in Bencraft/Joly files.]}
}
{
}

\Verse{79}
{
    \zR{只打听不出起程的日期，心里又是烦躁。}
}
{
    \eR{[English source not present in the checked-in Bencraft/Joly files.]}
}
{
}

\Verse{80}
{
    \zR{想到盘费 算来不敷，不得已写书一封，差人到赖尚荣任上借银五百，叫人沿途迎来，应付需 用。}
}
{
    \eR{[English source not present in the checked-in Bencraft/Joly files.]}
}
{
}

\Verse{81}
{
    \zR{过了数日，贾政的船才行得十数里，那家人回来，迎上船只，将赖尚荣的禀启 呈上。}
}
{
    \eR{[English source not present in the checked-in Bencraft/Joly files.]}
}
{
}

\Verse{82}
{
    \zR{书内告了多少苦处，备上白银五十两。}
}
{
    \eR{[English source not present in the checked-in Bencraft/Joly files.]}
}
{
}

\Verse{83}
{
    \zR{贾政看了大怒，即命家人：“立刻送 还!将原书发回，叫他不必费心。”}
}
{
    \eR{[English source not present in the checked-in Bencraft/Joly files.]}
}
{
}

\Verse{84}
{
    \zR{那家人无奈，只得回到赖尚荣任所。}
}
{
    \eR{[English source not present in the checked-in Bencraft/Joly files.]}
}
{
}

\Verse{85}
{
    \zR{赖尚荣接到 原书银两，心中烦闷，知事办得不周到，又添了一百，央来人带回，帮着说些好话。}
}
{
    \eR{[English source not present in the checked-in Bencraft/Joly files.]}
}
{
}

\Verse{86}
{
    \zR{岂知那人不肯带回，撂下就走。}
}
{
    \eR{[English source not present in the checked-in Bencraft/Joly files.]}
}
{
}

\Verse{87}
{
    \zR{赖尚荣心下不安，立刻修书到家，回明他父亲，叫 他设法告假，赎出身来。}
}
{
    \eR{[English source not present in the checked-in Bencraft/Joly files.]}
}
{
}

\Verse{88}
{
    \zR{于是赖家托了贾蔷贾芸等在王夫人面前乞恩放出。}
}
{
    \eR{[English source not present in the checked-in Bencraft/Joly files.]}
}
{
}

\Verse{89}
{
    \zR{贾蔷明 知不能，过了一日，假说王夫人不依的话，回覆了。}
}
{
    \eR{[English source not present in the checked-in Bencraft/Joly files.]}
}
{
}

\Verse{90}
{
    \zR{赖家一面告假，一面差人到赖 尚荣任上，叫他告病辞官。}
}
{
    \eR{[English source not present in the checked-in Bencraft/Joly files.]}
}
{
}

\Verse{91}
{
    \zR{王夫人并不知道。}
}
{
    \eR{[English source not present in the checked-in Bencraft/Joly files.]}
}
{
}

\Verse{92}
{
    \zR{那贾芸听见贾蔷的假话，心里便没想头。}
}
{
    \eR{[English source not present in the checked-in Bencraft/Joly files.]}
}
{
}

\Verse{93}
{
    \zR{连日在外又输了好些银钱，无所抵偿， 便和贾环借贷。}
}
{
    \eR{[English source not present in the checked-in Bencraft/Joly files.]}
}
{
}

\Verse{94}
{
    \zR{贾环本是一个钱没有的，虽是赵姨娘有些积蓄，早被他弄光了，那 能照应人家？}
}
{
    \eR{[English source not present in the checked-in Bencraft/Joly files.]}
}
{
}

\Verse{95}
{
    \zR{便想起凤姐待他刻薄，趁着贾琏不在家，要摆布巧姐出气，遂把这个 当叫贾芸来上，故意的埋怨贾芸道：“你们年纪又大，放着弄银钱的事又不敢办，
        倒和我没有钱的人商量。”}
}
{
    \eR{[English source not present in the checked-in Bencraft/Joly files.]}
}
{
}

\Verse{96}
{
    \zR{贾芸道：“三叔你这话说的倒好笑。}
}
{
    \eR{[English source not present in the checked-in Bencraft/Joly files.]}
}
{
}

\Verse{97}
{
    \zR{咱们一块儿玩，一块 儿闹，那里有有钱的事？”}
}
{
    \eR{[English source not present in the checked-in Bencraft/Joly files.]}
}
{
}

\Verse{98}
{
    \zR{贾环道：“不是前儿有人说是外藩要买个偏房？}
}
{
    \eR{[English source not present in the checked-in Bencraft/Joly files.]}
}
{
}

\Verse{99}
{
    \zR{你们何不 和王大舅商量，把巧姐说给他呢？”}
}
{
    \eR{[English source not present in the checked-in Bencraft/Joly files.]}
}
{
}

\Verse{100}
{
    \zR{贾芸道：“叔叔，我说句招你生气的话：外藩 花了钱买人，还想能和咱们走动么？”}
}
{
    \eR{[English source not present in the checked-in Bencraft/Joly files.]}
}
{
}

\Verse{101}
{
    \zR{贾环在贾芸耳边说了些话。}
}
{
    \eR{[English source not present in the checked-in Bencraft/Joly files.]}
}
{
}

\Verse{102}
{
    \zR{贾芸虽然点头， 只道贾环是小孩子的话，也不当事。}
}
{
    \eR{[English source not present in the checked-in Bencraft/Joly files.]}
}
{
}

\Verse{103}
{
    \zR{恰好王仁走来说道：“你们两个人商量些什么？}
}
{
    \eR{[English source not present in the checked-in Bencraft/Joly files.]}
}
{
}

\Verse{104}
{
    \zR{瞒着我吗？”}
}
{
    \eR{[English source not present in the checked-in Bencraft/Joly files.]}
}
{
}

\Verse{105}
{
    \zR{贾芸便将贾环的话附耳低言的说了。}
}
{
    \eR{[English source not present in the checked-in Bencraft/Joly files.]}
}
{
}

\Verse{106}
{
    \zR{王仁拍手道：“这倒是一宗好事， 又有银子。}
}
{
    \eR{[English source not present in the checked-in Bencraft/Joly files.]}
}
{
}

\Verse{107}
{
    \zR{只怕你们不能。}
}
{
    \eR{[English source not present in the checked-in Bencraft/Joly files.]}
}
{
}

\Verse{108}
{
    \zR{若是你们敢办，我是亲舅舅，做得主的。}
}
{
    \eR{[English source not present in the checked-in Bencraft/Joly files.]}
}
{
}

\Verse{109}
{
    \zR{只要环老三在 大太太跟前那么一说，我找邢大舅再一说，太太们问起来，你们打伙儿说好就是了。”}
}
{
    \eR{[English source not present in the checked-in Bencraft/Joly files.]}
}
{
}

\Verse{110}
{
    \zR{贾环等商议定了，王仁便去找邢大舅，贾芸便去回邢王二夫人，说得锦上添花。}
}
{
    \eR{[English source not present in the checked-in Bencraft/Joly files.]}
}
{
}

\Verse{111}
{
    \zR{王夫人听了，虽然入耳，只是不信。}
}
{
    \eR{[English source not present in the checked-in Bencraft/Joly files.]}
}
{
}

\Verse{112}
{
    \zR{邢夫人听得邢大舅知道，心里愿意，便打发人 找了邢大舅来问他。}
}
{
    \eR{[English source not present in the checked-in Bencraft/Joly files.]}
}
{
}

\Verse{113}
{
    \zR{那邢大舅已经听了王仁的话，又可分肥，便在邢夫人跟前说道： “若说这位郡王，是极有体面的。}
}
{
    \eR{[English source not present in the checked-in Bencraft/Joly files.]}
}
{
}

\Verse{114}
{
    \zR{若应了这门亲事，虽说不是正配，管保一过了门， 姐夫的官早复了，这里的声势又好了。”}
}
{
    \eR{[English source not present in the checked-in Bencraft/Joly files.]}
}
{
}

\Verse{115}
{
    \zR{邢夫人本是没主意的人，被傻大舅一番假 话哄得心动，请了王仁来一问，更说得热闹。}
}
{
    \eR{[English source not present in the checked-in Bencraft/Joly files.]}
}
{
}

\Verse{116}
{
    \zR{于是邢夫人倒叫人出去追着贾芸去说。}
}
{
    \eR{[English source not present in the checked-in Bencraft/Joly files.]}
}
{
}

\Verse{117}
{
    \zR{王仁即刻找了人去到外藩公馆说了。}
}
{
    \eR{[English source not present in the checked-in Bencraft/Joly files.]}
}
{
}

\Verse{118}
{
    \zR{那外藩不知底细，便要打发人来相看。}
}
{
    \eR{[English source not present in the checked-in Bencraft/Joly files.]}
}
{
}

\Verse{119}
{
    \zR{贾芸又 钻了相看的人，说明：“原是瞒着合宅的，只说是王府相亲。}
}
{
    \eR{[English source not present in the checked-in Bencraft/Joly files.]}
}
{
}

\Verse{120}
{
    \zR{等到成了，他祖母作 主，亲舅舅的保山，是不怕的。”}
}
{
    \eR{[English source not present in the checked-in Bencraft/Joly files.]}
}
{
}

\Verse{121}
{
    \zR{那相看的人应了。}
}
{
    \eR{[English source not present in the checked-in Bencraft/Joly files.]}
}
{
}

\Verse{122}
{
    \zR{贾芸便送信与邢夫人，并回了 王夫人。}
}
{
    \eR{[English source not present in the checked-in Bencraft/Joly files.]}
}
{
}

\Verse{123}
{
    \zR{那李纨宝钗等不知原故，只道是件好事，也都欢喜。}
}
{
    \eR{[English source not present in the checked-in Bencraft/Joly files.]}
}
{
}

\Verse{124}
{
    \zR{那日果然来了几个女人，都是艳妆丽服。}
}
{
    \eR{[English source not present in the checked-in Bencraft/Joly files.]}
}
{
}

\Verse{125}
{
    \zR{邢夫人接了进去，叙了些闲话。}
}
{
    \eR{[English source not present in the checked-in Bencraft/Joly files.]}
}
{
}

\Verse{126}
{
    \zR{那来 人本知是个诰命，也不敢怠慢。}
}
{
    \eR{[English source not present in the checked-in Bencraft/Joly files.]}
}
{
}

\Verse{127}
{
    \zR{邢夫人因事未定，也没有和巧姐说明，只说有亲戚 来瞧，叫他去见。}
}
{
    \eR{[English source not present in the checked-in Bencraft/Joly files.]}
}
{
}

\Verse{128}
{
    \zR{巧姐到底是个小孩子，那管这些，便跟了奶妈过来。}
}
{
    \eR{[English source not present in the checked-in Bencraft/Joly files.]}
}
{
}

\Verse{129}
{
    \zR{平儿不放心， 也跟着来。}
}
{
    \eR{[English source not present in the checked-in Bencraft/Joly files.]}
}
{
}

\Verse{130}
{
    \zR{只见有两个宫人打扮的，见了巧姐，便浑身上下一看，更又起身来拉着 巧姐的手又瞧了一遍，略坐了一坐就走了。}
}
{
    \eR{[English source not present in the checked-in Bencraft/Joly files.]}
}
{
}

\Verse{131}
{
    \zR{倒把巧姐看得羞臊。}
}
{
    \eR{[English source not present in the checked-in Bencraft/Joly files.]}
}
{
}

\Verse{132}
{
    \zR{回到房中纳闷，想 来没有这门亲戚，便问平儿。}
}
{
    \eR{[English source not present in the checked-in Bencraft/Joly files.]}
}
{
}

\Verse{133}
{
    \zR{平儿先看见来头，却也猜着八九：“必是相亲的。}
}
{
    \eR{[English source not present in the checked-in Bencraft/Joly files.]}
}
{
}

\Verse{134}
{
    \zR{但 是二爷不在家，大太太作主，到底不知是那府里的。}
}
{
    \eR{[English source not present in the checked-in Bencraft/Joly files.]}
}
{
}

\Verse{135}
{
    \zR{若说是对头亲，不该这样相看。}
}
{
    \eR{[English source not present in the checked-in Bencraft/Joly files.]}
}
{
}

\Verse{136}
{
    \zR{瞧那几个人的来头，不像是本支王府，好像是外头路数。}
}
{
    \eR{[English source not present in the checked-in Bencraft/Joly files.]}
}
{
}

\Verse{137}
{
    \zR{如今且不必和姑娘说明， 且打听明白再说。”}
}
{
    \eR{[English source not present in the checked-in Bencraft/Joly files.]}
}
{
}

\Verse{138}
{
    \zR{平儿心下留神打听，那些丫头婆子都是平儿使过的，平儿一问，所有听见外头 的风声都告诉了。}
}
{
    \eR{[English source not present in the checked-in Bencraft/Joly files.]}
}
{
}

\Verse{139}
{
    \zR{平儿便吓的没了主意，虽不和巧姐说，便赶着去告诉了李纨宝钗， 求他二人告诉王夫人。}
}
{
    \eR{[English source not present in the checked-in Bencraft/Joly files.]}
}
{
}

\Verse{140}
{
    \zR{王夫人知道这事不好，便和邢夫人说知。}
}
{
    \eR{[English source not present in the checked-in Bencraft/Joly files.]}
}
{
}

\Verse{141}
{
    \zR{怎奈邢夫人信了兄 弟并王仁的话，反疑心王夫人不是好意，便说：“孙女儿也大了。}
}
{
    \eR{[English source not present in the checked-in Bencraft/Joly files.]}
}
{
}

\Verse{142}
{
    \zR{现在琏儿不在家， 这件事我还做得主。}
}
{
    \eR{[English source not present in the checked-in Bencraft/Joly files.]}
}
{
}

\Verse{143}
{
    \zR{况且他亲舅爷爷和他亲舅舅打听的，难道倒比别人不真么？}
}
{
    \eR{[English source not present in the checked-in Bencraft/Joly files.]}
}
{
}

\Verse{144}
{
    \zR{我 横竖是愿意的。}
}
{
    \eR{[English source not present in the checked-in Bencraft/Joly files.]}
}
{
}

\Verse{145}
{
    \zR{倘有什么不好，我和琏儿也抱怨不着别人。”}
}
{
    \eR{[English source not present in the checked-in Bencraft/Joly files.]}
}
{
}

\Verse{146}
{
    \zR{王夫人听了这些话， 心下暗暗生气，勉强说些闲话，便走了出来告诉了宝钗，自己落泪。}
}
{
    \eR{[English source not present in the checked-in Bencraft/Joly files.]}
}
{
}

\Verse{147}
{
    \zR{宝玉劝道：“太 太别烦恼。}
}
{
    \eR{[English source not present in the checked-in Bencraft/Joly files.]}
}
{
}

\Verse{148}
{
    \zR{这件事，我看来是不成的。}
}
{
    \eR{[English source not present in the checked-in Bencraft/Joly files.]}
}
{
}

\Verse{149}
{
    \zR{这又是巧姐儿命里所招，只求太太不管就是 了。”}
}
{
    \eR{[English source not present in the checked-in Bencraft/Joly files.]}
}
{
}

\Verse{150}
{
    \zR{王夫人道：“你一开口就是疯话!人家说定了就要接过去。}
}
{
    \eR{[English source not present in the checked-in Bencraft/Joly files.]}
}
{
}

\Verse{151}
{
    \zR{若依平儿的话，你 琏二哥哥不抱怨我么？}
}
{
    \eR{[English source not present in the checked-in Bencraft/Joly files.]}
}
{
}

\Verse{152}
{
    \zR{别说自己的侄孙女儿，就是亲戚家的，也是要好才好。}
}
{
    \eR{[English source not present in the checked-in Bencraft/Joly files.]}
}
{
}

\Verse{153}
{
    \zR{邢姑 娘是我们作媒的，配了你二大舅子，如今和和顺顺的过日子，不好么？}
}
{
    \eR{[English source not present in the checked-in Bencraft/Joly files.]}
}
{
}

\Verse{154}
{
    \zR{那琴姑娘， 梅家娶了去，听见说是丰衣足食的，很好。}
}
{
    \eR{[English source not present in the checked-in Bencraft/Joly files.]}
}
{
}

\Verse{155}
{
    \zR{就是史姑娘，是他叔叔的主意，头里原 好，如今姑爷痨病死了，你史妹妹立志守寡，也就苦了。}
}
{
    \eR{[English source not present in the checked-in Bencraft/Joly files.]}
}
{
}

\Verse{156}
{
    \zR{若是巧姐儿错给了人家儿， 可不是我的心坏？”}
}
{
    \eR{[English source not present in the checked-in Bencraft/Joly files.]}
}
{
}

\Verse{157}
{
    \zR{正说着，平儿过来瞧宝钗，并探听邢夫人的口气。}
}
{
    \eR{[English source not present in the checked-in Bencraft/Joly files.]}
}
{
}

\Verse{158}
{
    \zR{王夫人将邢 夫人的话说了一遍。}
}
{
    \eR{[English source not present in the checked-in Bencraft/Joly files.]}
}
{
}

\Verse{159}
{
    \zR{平儿呆了半天，跪下求道：“巧姐儿终身，全仗着太太!若信了 人家的话，不但姑娘一辈子受了苦，便是琏二爷回来，怎么说呢？”}
}
{
    \eR{[English source not present in the checked-in Bencraft/Joly files.]}
}
{
}

\Verse{160}
{
    \zR{王夫人道：“你 是个明白人，起来听我说：巧姐儿到底是大太太孙女儿，他要作主，我能够拦他么？”}
}
{
    \eR{[English source not present in the checked-in Bencraft/Joly files.]}
}
{
}

\Verse{161}
{
    \zR{宝玉劝道：“无妨碍的，只要明白就是了。”}
}
{
    \eR{[English source not present in the checked-in Bencraft/Joly files.]}
}
{
}

\Verse{162}
{
    \zR{平儿生怕宝玉疯癫嚷出来，也并不言语， 回了王夫人，竟自去了。}
}
{
    \eR{[English source not present in the checked-in Bencraft/Joly files.]}
}
{
}

\Verse{163}
{
    \zR{这里王夫人想到烦闷，一阵心痛，叫丫头扶着，勉强回到自己房中躺下，不叫 宝玉宝钗过来，说睡睡就好的。}
}
{
    \eR{[English source not present in the checked-in Bencraft/Joly files.]}
}
{
}

\Verse{164}
{
    \zR{自己却也烦闷。}
}
{
    \eR{[English source not present in the checked-in Bencraft/Joly files.]}
}
{
}

\Verse{165}
{
    \zR{听见说李婶娘来了，也不及接待。}
}
{
    \eR{[English source not present in the checked-in Bencraft/Joly files.]}
}
{
}

\Verse{166}
{
    \zR{只见贾兰进来请了安，回道：“今早爷爷那里打发人带了一封书子来，外头小子们 传进来的。}
}
{
    \eR{[English source not present in the checked-in Bencraft/Joly files.]}
}
{
}

\Verse{167}
{
    \zR{我母亲接了，正要过来，因我老娘来了，叫我先呈给太太瞧，回来我母 亲就过来来回太太。}
}
{
    \eR{[English source not present in the checked-in Bencraft/Joly files.]}
}
{
}

\Verse{168}
{
    \zR{还说我老娘要过来呢。”}
}
{
    \eR{[English source not present in the checked-in Bencraft/Joly files.]}
}
{
}

\Verse{169}
{
    \zR{说着，一面把书子呈上。}
}
{
    \eR{[English source not present in the checked-in Bencraft/Joly files.]}
}
{
}

\Verse{170}
{
    \zR{王夫人一面 接书，一面问道：“你老娘来作什么？”}
}
{
    \eR{[English source not present in the checked-in Bencraft/Joly files.]}
}
{
}

\Verse{171}
{
    \zR{贾兰道：“我也不知道。}
}
{
    \eR{[English source not present in the checked-in Bencraft/Joly files.]}
}
{
}

\Verse{172}
{
    \zR{我只听见我老娘说： 我三姨儿的婆婆家有什么信儿来了。”}
}
{
    \eR{[English source not present in the checked-in Bencraft/Joly files.]}
}
{
}

\Verse{173}
{
    \zR{王夫人听了，想起来还是前次给甄宝玉说了 李绮，后来放定下茶，想来此时甄家要娶过门，所以李婶娘来商量这件事情。}
}
{
    \eR{[English source not present in the checked-in Bencraft/Joly files.]}
}
{
}

\Verse{174}
{
    \zR{便点 点头儿，一面拆开书信，见上面写着道： 近因沿途俱系海疆凯旋船只，不能迅速前行。}
}
{
    \eR{[English source not present in the checked-in Bencraft/Joly files.]}
}
{
}

\Verse{175}
{
    \zR{闻探姐随翁婿来都，不知曾有信 否？}
}
{
    \eR{[English source not present in the checked-in Bencraft/Joly files.]}
}
{
}

\Verse{176}
{
    \zR{前接到琏侄手禀，知大老爷身体欠安，亦不知已有确信否？}
}
{
    \eR{[English source not present in the checked-in Bencraft/Joly files.]}
}
{
}

\Verse{177}
{
    \zR{宝玉兰儿场期已近， 务须实心用功，不可怠惰。}
}
{
    \eR{[English source not present in the checked-in Bencraft/Joly files.]}
}
{
}

\Verse{178}
{
    \zR{老太太灵柩抵家，尚需日时。}
}
{
    \eR{[English source not present in the checked-in Bencraft/Joly files.]}
}
{
}

\Verse{179}
{
    \zR{我身体平善，不必挂念。}
}
{
    \eR{[English source not present in the checked-in Bencraft/Joly files.]}
}
{
}

\Verse{180}
{
    \zR{此谕宝玉等知道。}
}
{
    \eR{[English source not present in the checked-in Bencraft/Joly files.]}
}
{
}

\Verse{181}
{
    \zR{月日手书。}
}
{
    \eR{[English source not present in the checked-in Bencraft/Joly files.]}
}
{
}

\Verse{182}
{
    \zR{蓉儿另禀。}
}
{
    \eR{[English source not present in the checked-in Bencraft/Joly files.]}
}
{
}

\Verse{183}
{
    \zR{王夫人看了，仍旧递给贾兰，说：“你拿去给你二叔叔瞧瞧，还交给你母亲罢。”}
}
{
    \eR{[English source not present in the checked-in Bencraft/Joly files.]}
}
{
}

\Verse{184}
{
    \zR{正 说着，李纨同李婶娘过来，请安问好毕，王夫人让了坐。}
}
{
    \eR{[English source not present in the checked-in Bencraft/Joly files.]}
}
{
}

\Verse{185}
{
    \zR{李婶娘便将甄家要娶李绮 的话说了一遍。}
}
{
    \eR{[English source not present in the checked-in Bencraft/Joly files.]}
}
{
}

\Verse{186}
{
    \zR{大家商议了一会子。}
}
{
    \eR{[English source not present in the checked-in Bencraft/Joly files.]}
}
{
}

\Verse{187}
{
    \zR{李纨因问王夫人道：“老爷的书子，太太看过 了么？”}
}
{
    \eR{[English source not present in the checked-in Bencraft/Joly files.]}
}
{
}

\Verse{188}
{
    \zR{王夫人道：“看过了。”}
}
{
    \eR{[English source not present in the checked-in Bencraft/Joly files.]}
}
{
}

\Verse{189}
{
    \zR{贾兰便拿着给他母亲瞧。}
}
{
    \eR{[English source not present in the checked-in Bencraft/Joly files.]}
}
{
}

\Verse{190}
{
    \zR{李纨看了道：“三姑娘出 了门好几年，总没有来，如今要回京了，太太也放了好些心。”}
}
{
    \eR{[English source not present in the checked-in Bencraft/Joly files.]}
}
{
}

\Verse{191}
{
    \zR{王夫人道：“我本是 心痛，看见探丫头要回来了，心里略好些，只是不知几时才到？”}
}
{
    \eR{[English source not present in the checked-in Bencraft/Joly files.]}
}
{
}

\Verse{192}
{
    \zR{李婶娘便问了贾 政在路好。}
}
{
    \eR{[English source not present in the checked-in Bencraft/Joly files.]}
}
{
}

\Verse{193}
{
    \zR{李纨因向贾兰道：“哥儿瞧见了？}
}
{
    \eR{[English source not present in the checked-in Bencraft/Joly files.]}
}
{
}

\Verse{194}
{
    \zR{场期近了，你爷爷惦记的什么似的。}
}
{
    \eR{[English source not present in the checked-in Bencraft/Joly files.]}
}
{
}

\Verse{195}
{
    \zR{你 快拿了去给二叔叔瞧去罢。”}
}
{
    \eR{[English source not present in the checked-in Bencraft/Joly files.]}
}
{
}

\Verse{196}
{
    \zR{李婶娘道：“他们爷儿两个又没进过学，怎么能下场 呢？”}
}
{
    \eR{[English source not present in the checked-in Bencraft/Joly files.]}
}
{
}

\Verse{197}
{
    \zR{王夫人道：“他爷爷做粮道的起身时，给他们爷儿两个援了例监了。”}
}
{
    \eR{[English source not present in the checked-in Bencraft/Joly files.]}
}
{
}

\Verse{198}
{
    \zR{李婶娘 点头。}
}
{
    \eR{[English source not present in the checked-in Bencraft/Joly files.]}
}
{
}

\Verse{199}
{
    \zR{贾兰一面拿着书子出来，来找宝玉。}
}
{
    \eR{[English source not present in the checked-in Bencraft/Joly files.]}
}
{
}

\Verse{200}
{
    \zR{却说宝玉送了王夫人去后，正拿着《秋水》一篇在那里细玩。}
}
{
    \eR{[English source not present in the checked-in Bencraft/Joly files.]}
}
{
}

\Verse{201}
{
    \zR{宝钗从里间走出， 见他看的得意忘言，便走过来一看。}
}
{
    \eR{[English source not present in the checked-in Bencraft/Joly files.]}
}
{
}

\Verse{202}
{
    \zR{见是这个，心里着实烦闷，细想：“他只顾把 这些出世离群的话当作一件正经事，终久不妥！”}
}
{
    \eR{[English source not present in the checked-in Bencraft/Joly files.]}
}
{
}

\Verse{203}
{
    \zR{看他这种光景，料劝不过来，便 坐在宝玉傍边，怔怔的瞅着。}
}
{
    \eR{[English source not present in the checked-in Bencraft/Joly files.]}
}
{
}

\Verse{204}
{
    \zR{宝玉见他这般，便道：“你这又是为什么？”}
}
{
    \eR{[English source not present in the checked-in Bencraft/Joly files.]}
}
{
}

\Verse{205}
{
    \zR{宝钗道： “我想你我既为夫妇，你便是我终身的倚靠，却不在情欲之私。}
}
{
    \eR{[English source not present in the checked-in Bencraft/Joly files.]}
}
{
}

\Verse{206}
{
    \zR{论起荣华富贵，原 不过是过眼烟云；但自古圣贤，以人品根柢为重――”宝玉也没听完，把那本书搁
        在旁边，微微的笑道：“据你说‘人品根柢’，又是什么‘古圣贤’，你可知古圣贤 说过，‘不失其赤子之心’？}
}
{
    \eR{[English source not present in the checked-in Bencraft/Joly files.]}
}
{
}

\Verse{207}
{
    \zR{那赤子有什么好处？}
}
{
    \eR{[English source not present in the checked-in Bencraft/Joly files.]}
}
{
}

\Verse{208}
{
    \zR{不过是无知无识无贪无忌。}
}
{
    \eR{[English source not present in the checked-in Bencraft/Joly files.]}
}
{
}

\Verse{209}
{
    \zR{我们生 来已陷溺在贪嗔痴爱中，犹如污泥一般，怎么能跳出这般尘网？}
}
{
    \eR{[English source not present in the checked-in Bencraft/Joly files.]}
}
{
}

\Verse{210}
{
    \zR{如今才晓得‘聚散 浮生’四字，古人说了，不曾提醒一个。}
}
{
    \eR{[English source not present in the checked-in Bencraft/Joly files.]}
}
{
}

\Verse{211}
{
    \zR{既要讲到人品根柢，谁是到那太初一步地 位的？”}
}
{
    \eR{[English source not present in the checked-in Bencraft/Joly files.]}
}
{
}

\Verse{212}
{
    \zR{宝钗道：“你既说‘赤子之心’，古圣贤原以忠孝为赤子之心，并不是遁世 离群、无关无系为赤子之心。}
}
{
    \eR{[English source not present in the checked-in Bencraft/Joly files.]}
}
{
}

\Verse{213}
{
    \zR{尧、舜、禹、汤、周、孔，时刻以救民济世为心，所 谓赤子之心，原不过是‘不忍’二字。}
}
{
    \eR{[English source not present in the checked-in Bencraft/Joly files.]}
}
{
}

\Verse{214}
{
    \zR{若你方才所说的忍于抛弃天伦，还成什么道 理？”}
}
{
    \eR{[English source not present in the checked-in Bencraft/Joly files.]}
}
{
}

\Verse{215}
{
    \zR{宝玉点头笑道：“尧舜不强巢许，武周不强夷齐。”}
}
{
    \eR{[English source not present in the checked-in Bencraft/Joly files.]}
}
{
}

\Verse{216}
{
    \zR{宝钗不等他说完，便道： “你这个话，益发不是了。}
}
{
    \eR{[English source not present in the checked-in Bencraft/Joly files.]}
}
{
}

\Verse{217}
{
    \zR{古来若都是巢、许、夷、齐，为什么如今人又把尧、舜、 周、孔称为圣贤呢？}
}
{
    \eR{[English source not present in the checked-in Bencraft/Joly files.]}
}
{
}

\Verse{218}
{
    \zR{况且你自比夷齐，更不成话。}
}
{
    \eR{[English source not present in the checked-in Bencraft/Joly files.]}
}
{
}

\Verse{219}
{
    \zR{夷齐原是生在殷商末世，有许多 难处之事，所以才有托而逃。}
}
{
    \eR{[English source not present in the checked-in Bencraft/Joly files.]}
}
{
}

\Verse{220}
{
    \zR{当此圣世，咱们世受国恩，祖父锦衣玉食；况你自有 生以来，自去世的老太太，以及老爷太太，视如珍宝。}
}
{
    \eR{[English source not present in the checked-in Bencraft/Joly files.]}
}
{
}

\Verse{221}
{
    \zR{你方才所说，自己想一想， 是与不是？”}
}
{
    \eR{[English source not present in the checked-in Bencraft/Joly files.]}
}
{
}

\Verse{222}
{
    \zR{宝玉听了，也不答言，只有仰头微笑。}
}
{
    \eR{[English source not present in the checked-in Bencraft/Joly files.]}
}
{
}

\Verse{223}
{
    \zR{宝钗因又劝道：“你既理屈词 穷，我劝你从此把心收一收，好好的用用功，但能博得一第，便是从此而止，也不 枉天恩祖德了。”}
}
{
    \eR{[English source not present in the checked-in Bencraft/Joly files.]}
}
{
}

\Verse{224}
{
    \zR{宝玉点了点头，叹了口气，说道：“一第呢其实也不是什么难事。}
}
{
    \eR{[English source not present in the checked-in Bencraft/Joly files.]}
}
{
}

\Verse{225}
{
    \zR{倒是你这个‘从此而止’，‘不枉天恩祖德’，却还不离其宗。”}
}
{
    \eR{[English source not present in the checked-in Bencraft/Joly files.]}
}
{
}

\Verse{226}
{
    \zR{宝钗未及答言，袭人 过来说道：“刚才二奶奶说的古圣先贤，我们也不懂。}
}
{
    \eR{[English source not present in the checked-in Bencraft/Joly files.]}
}
{
}

\Verse{227}
{
    \zR{我只想着我们这些人，从小 儿辛辛苦苦跟着二爷，不知陪了多少小心，论起理来原该当的，但只二爷也该体谅 体谅。}
}
{
    \eR{[English source not present in the checked-in Bencraft/Joly files.]}
}
{
}

\Verse{228}
{
    \zR{况且二奶奶替二爷在老爷太太跟前行了多少孝道，就是二爷不以夫妻为事， 也不可太辜负了人心。}
}
{
    \eR{[English source not present in the checked-in Bencraft/Joly files.]}
}
{
}

\Verse{229}
{
    \zR{至于神仙那一层，更是谎话，谁见过有走到凡间来的神仙呢？}
}
{
    \eR{[English source not present in the checked-in Bencraft/Joly files.]}
}
{
}

\Verse{230}
{
    \zR{那里来的这么个和尚，说了些混话，二爷就信了真!二爷是读书的人，难道他的话 比老爷太太还重么？”}
}
{
    \eR{[English source not present in the checked-in Bencraft/Joly files.]}
}
{
}

\Verse{231}
{
    \zR{宝玉听了，低头不语。}
}
{
    \eR{[English source not present in the checked-in Bencraft/Joly files.]}
}
{
}

\Verse{232}
{
    \zR{袭人还要说时，只听外面脚步走响，隔着窗户问道：“二叔在屋里呢么？”}
}
{
    \eR{[English source not present in the checked-in Bencraft/Joly files.]}
}
{
}

\Verse{233}
{
    \zR{宝 玉听了是贾兰的声音，便站起来笑道：“你进来罢。”}
}
{
    \eR{[English source not present in the checked-in Bencraft/Joly files.]}
}
{
}

\Verse{234}
{
    \zR{宝钗也站起来。}
}
{
    \eR{[English source not present in the checked-in Bencraft/Joly files.]}
}
{
}

\Verse{235}
{
    \zR{贾兰进来，笑 容可掬的给宝玉宝钗请了安，问了袭人的好，袭人也问了好，便把书子呈给宝玉瞧。}
}
{
    \eR{[English source not present in the checked-in Bencraft/Joly files.]}
}
{
}

\Verse{236}
{
    \zR{宝玉接在手中看了，便道：“你三姑姑回来了？”}
}
{
    \eR{[English source not present in the checked-in Bencraft/Joly files.]}
}
{
}

\Verse{237}
{
    \zR{贾兰道：“爷爷既如此写，自然是 回来的了。”}
}
{
    \eR{[English source not present in the checked-in Bencraft/Joly files.]}
}
{
}

\Verse{238}
{
    \zR{宝玉点头不语，默默如有所思。}
}
{
    \eR{[English source not present in the checked-in Bencraft/Joly files.]}
}
{
}

\Verse{239}
{
    \zR{贾兰便问：“叔叔看见了：爷爷后头写 着，叫咱们好生念书呢。}
}
{
    \eR{[English source not present in the checked-in Bencraft/Joly files.]}
}
{
}

\Verse{240}
{
    \zR{叔叔这成子只怕总没作文章罢？”}
}
{
    \eR{[English source not present in the checked-in Bencraft/Joly files.]}
}
{
}

\Verse{241}
{
    \zR{宝玉笑道：“我也要作 几篇熟一熟手，好去诓这个功名。”}
}
{
    \eR{[English source not present in the checked-in Bencraft/Joly files.]}
}
{
}

\Verse{242}
{
    \zR{贾兰道：“叔叔既这样，就拟几个题目，我跟着 叔叔作作，也好进去混场。}
}
{
    \eR{[English source not present in the checked-in Bencraft/Joly files.]}
}
{
}

\Verse{243}
{
    \zR{别到那时交了白卷子，惹人笑话；不但笑话我，人家连 叔叔都要笑话了。”}
}
{
    \eR{[English source not present in the checked-in Bencraft/Joly files.]}
}
{
}

\Verse{244}
{
    \zR{宝玉道：“你也不至如此。”}
}
{
    \eR{[English source not present in the checked-in Bencraft/Joly files.]}
}
{
}

\Verse{245}
{
    \zR{说着，宝钗命贾兰坐下。}
}
{
    \eR{[English source not present in the checked-in Bencraft/Joly files.]}
}
{
}

\Verse{246}
{
    \zR{宝玉仍坐 在原处，贾兰侧身坐了。}
}
{
    \eR{[English source not present in the checked-in Bencraft/Joly files.]}
}
{
}

\Verse{247}
{
    \zR{两个谈了一回文，不觉喜动颜色。}
}
{
    \eR{[English source not present in the checked-in Bencraft/Joly files.]}
}
{
}

\Verse{248}
{
    \zR{宝钗见他爷儿两个谈得 高兴，便仍进屋里去了，心中细想：“宝玉此时光景，或者醒悟过来了。}
}
{
    \eR{[English source not present in the checked-in Bencraft/Joly files.]}
}
{
}

\Verse{249}
{
    \zR{只是刚才 说话，他把那‘从此而止’四字单单的许可，这又不知是什么意思了？”}
}
{
    \eR{[English source not present in the checked-in Bencraft/Joly files.]}
}
{
}

\Verse{250}
{
    \zR{宝钗尚自 犹豫。}
}
{
    \eR{[English source not present in the checked-in Bencraft/Joly files.]}
}
{
}

\Verse{251}
{
    \zR{惟有袭人看他爱讲文章，提到下场，更又欣然，心里想道：“阿弥陀佛!好容 易讲《四书》似的才讲过来了。”}
}
{
    \eR{[English source not present in the checked-in Bencraft/Joly files.]}
}
{
}

\Verse{252}
{
    \zR{这里宝玉和贾兰讲文，莺儿沏过茶来。}
}
{
    \eR{[English source not present in the checked-in Bencraft/Joly files.]}
}
{
}

\Verse{253}
{
    \zR{贾兰站起 来接了，又说了一会子下场的规矩，并请甄宝玉在一处的话，宝玉也甚似愿意。}
}
{
    \eR{[English source not present in the checked-in Bencraft/Joly files.]}
}
{
}

\Verse{254}
{
    \zR{一时贾兰回去，便将书子留给宝玉了。}
}
{
    \eR{[English source not present in the checked-in Bencraft/Joly files.]}
}
{
}

\Verse{255}
{
    \zR{那宝玉看着书子，笑嘻嘻走进来，递给 麝月收了，便出来将那本《庄子》收了。}
}
{
    \eR{[English source not present in the checked-in Bencraft/Joly files.]}
}
{
}

\Verse{256}
{
    \zR{把几部向来最得意的，如《参同契》、《元 命苞》、《五灯会元》之类，叫出麝月、秋纹、莺儿等都搬了搁在一边。}
}
{
    \eR{[English source not present in the checked-in Bencraft/Joly files.]}
}
{
}

\Verse{257}
{
    \zR{宝钗见他这 番举动，甚为罕异，因欲试探他，便笑问道：“不看他倒是正经，但又何必搬开呢。”}
}
{
    \eR{[English source not present in the checked-in Bencraft/Joly files.]}
}
{
}

\Verse{258}
{
    \zR{宝玉道：“如今才明白过来了。}
}
{
    \eR{[English source not present in the checked-in Bencraft/Joly files.]}
}
{
}

\Verse{259}
{
    \zR{这些书都算不得什么。}
}
{
    \eR{[English source not present in the checked-in Bencraft/Joly files.]}
}
{
}

\Verse{260}
{
    \zR{我还要一火焚之，方为干净。”}
}
{
    \eR{[English source not present in the checked-in Bencraft/Joly files.]}
}
{
}

\Verse{261}
{
    \zR{宝钗听了，更欣喜异常。}
}
{
    \eR{[English source not present in the checked-in Bencraft/Joly files.]}
}
{
}

\Verse{262}
{
    \zR{只听宝玉口中微吟道：“内典语中无佛性，金丹法外有仙 舟。”}
}
{
    \eR{[English source not present in the checked-in Bencraft/Joly files.]}
}
{
}

\Verse{263}
{
    \zR{宝钗也没很听真，只听得“无佛性”，“有仙舟”几个字，心中转又狐疑，且 看他作何光景。}
}
{
    \eR{[English source not present in the checked-in Bencraft/Joly files.]}
}
{
}

\Verse{264}
{
    \zR{宝玉便命麝月秋纹等收拾一间静室，把那些语录名稿及应制诗之类 都找出来，搁在静室中，自己却当真静静的用起功来。}
}
{
    \eR{[English source not present in the checked-in Bencraft/Joly files.]}
}
{
}

\Verse{265}
{
    \zR{宝钗这才放了心。}
}
{
    \eR{[English source not present in the checked-in Bencraft/Joly files.]}
}
{
}

\Verse{266}
{
    \zR{那袭人此时真是闻所未闻，见所未见，便悄悄的笑着向宝钗道：“到底奶奶说 话透彻!只一路讲究，就把二爷劝明白了。}
}
{
    \eR{[English source not present in the checked-in Bencraft/Joly files.]}
}
{
}

\Verse{267}
{
    \zR{就只可惜迟了一点儿，临场太近了。”}
}
{
    \eR{[English source not present in the checked-in Bencraft/Joly files.]}
}
{
}

\Verse{268}
{
    \zR{宝 钗点头微笑道：“功名自有定数，中与不中，倒也不在用功的迟早。}
}
{
    \eR{[English source not present in the checked-in Bencraft/Joly files.]}
}
{
}

\Verse{269}
{
    \zR{但愿他从此一 心巴结正路，把从前那些邪魔永不沾染，就是好了。”}
}
{
    \eR{[English source not present in the checked-in Bencraft/Joly files.]}
}
{
}

\Verse{270}
{
    \zR{说到这里，见房里无人，便 悄说道：“这一番悔悟过来固然很好，但只一件：怕又犯了前头的旧病，和女孩儿 们打起交道来，也是不好。”}
}
{
    \eR{[English source not present in the checked-in Bencraft/Joly files.]}
}
{
}

\Verse{271}
{
    \zR{袭人道：“奶奶说的也是。}
}
{
    \eR{[English source not present in the checked-in Bencraft/Joly files.]}
}
{
}

\Verse{272}
{
    \zR{二爷自从信了和尚，才把这 些姐妹冷淡了；如今不信和尚，真怕又要犯了前头的旧病呢。}
}
{
    \eR{[English source not present in the checked-in Bencraft/Joly files.]}
}
{
}

\Verse{273}
{
    \zR{我想：奶奶和我，二 爷原不大理会。}
}
{
    \eR{[English source not present in the checked-in Bencraft/Joly files.]}
}
{
}

\Verse{274}
{
    \zR{紫鹃去了，如今只他们四个。}
}
{
    \eR{[English source not present in the checked-in Bencraft/Joly files.]}
}
{
}

\Verse{275}
{
    \zR{这里头就是五儿有些个狐媚子，听见 说，他妈求了大奶奶和奶奶，说要讨出去给人家儿呢，但是这两天到底在这里呢。}
}
{
    \eR{[English source not present in the checked-in Bencraft/Joly files.]}
}
{
}

\Verse{276}
{
    \zR{麝月秋纹虽没别的，只是二爷那几年也都有些顽顽皮皮的。}
}
{
    \eR{[English source not present in the checked-in Bencraft/Joly files.]}
}
{
}

\Verse{277}
{
    \zR{如今算来，只有莺儿二 爷倒不大理会，况且莺儿也稳重。}
}
{
    \eR{[English source not present in the checked-in Bencraft/Joly files.]}
}
{
}

\Verse{278}
{
    \zR{我想倒茶弄水，只叫莺儿带着小丫头们伏侍就够 了，不知奶奶心里怎么样？”}
}
{
    \eR{[English source not present in the checked-in Bencraft/Joly files.]}
}
{
}

\Verse{279}
{
    \zR{宝钗道：“我也虑的是这个，你说的倒也罢了。”}
}
{
    \eR{[English source not present in the checked-in Bencraft/Joly files.]}
}
{
}

\Verse{280}
{
    \zR{从此 便派莺儿带着小丫头伏侍。}
}
{
    \eR{[English source not present in the checked-in Bencraft/Joly files.]}
}
{
}

\Verse{281}
{
    \zR{那宝玉却也不出房门，天天只差人去给王夫人请安。}
}
{
    \eR{[English source not present in the checked-in Bencraft/Joly files.]}
}
{
}

\Verse{282}
{
    \zR{王 夫人听见他这番光景，那一种欣慰之情更不待言了。}
}
{
    \eR{[English source not present in the checked-in Bencraft/Joly files.]}
}
{
}

\Verse{283}
{
    \zR{到了八月初三这一日，正是贾母的冥寿。}
}
{
    \eR{[English source not present in the checked-in Bencraft/Joly files.]}
}
{
}

\Verse{284}
{
    \zR{宝玉早晨过来磕了头，便回去，仍到 静室中去了。}
}
{
    \eR{[English source not present in the checked-in Bencraft/Joly files.]}
}
{
}

\Verse{285}
{
    \zR{饭后，宝钗袭人等都和姊妹们跟着邢王二夫人在前面屋里说闲话儿。}
}
{
    \eR{[English source not present in the checked-in Bencraft/Joly files.]}
}
{
}

\Verse{286}
{
    \zR{宝玉自在静室，冥心危坐。}
}
{
    \eR{[English source not present in the checked-in Bencraft/Joly files.]}
}
{
}

\Verse{287}
{
    \zR{忽见莺儿端了一盘瓜果进来，说：“太太叫人送来给二 爷吃的，这是老太太的克什。”}
}
{
    \eR{[English source not present in the checked-in Bencraft/Joly files.]}
}
{
}

\Verse{288}
{
    \zR{宝玉站起来答应了，复又坐下，便道：“搁在那里罢。”}
}
{
    \eR{[English source not present in the checked-in Bencraft/Joly files.]}
}
{
}

\Verse{289}
{
    \zR{莺儿一面放下瓜果，一面悄悄向宝玉道：“太太那里夸二爷呢。”}
}
{
    \eR{[English source not present in the checked-in Bencraft/Joly files.]}
}
{
}

\Verse{290}
{
    \zR{宝玉微笑。}
}
{
    \eR{[English source not present in the checked-in Bencraft/Joly files.]}
}
{
}

\Verse{291}
{
    \zR{莺儿又 道：“太太说了：二爷这一用功，明儿进场中了出来，明年再中了进士，作了官， 老爷太太可就不枉了盼二爷了。”}
}
{
    \eR{[English source not present in the checked-in Bencraft/Joly files.]}
}
{
}

\Verse{292}
{
    \zR{宝玉也只点头微笑。}
}
{
    \eR{[English source not present in the checked-in Bencraft/Joly files.]}
}
{
}

\Verse{293}
{
    \zR{莺儿忽然想起那年给宝玉打 络子的时候宝玉说的话来，便道：“真要二爷中了，那可是我们姑奶奶的造化了。}
}
{
    \eR{[English source not present in the checked-in Bencraft/Joly files.]}
}
{
}

\Verse{294}
{
    \zR{二爷还记得那一年在园子里，不是二爷叫我打梅花络子时说的：我们姑奶奶后来带 着我不知到那一个有造化的人家儿去呢？}
}
{
    \eR{[English source not present in the checked-in Bencraft/Joly files.]}
}
{
}

\Verse{295}
{
    \zR{如今二爷可是有造化的罢咧！”}
}
{
    \eR{[English source not present in the checked-in Bencraft/Joly files.]}
}
{
}

\Verse{296}
{
    \zR{宝玉听到这 里，又觉尘心一动，连忙敛神定息，微微的笑道：“据你说来，我是有造化的，你 们姑娘也是有造化的，你呢？”}
}
{
    \eR{[English source not present in the checked-in Bencraft/Joly files.]}
}
{
}

\Verse{297}
{
    \zR{莺儿把脸飞红了，勉强笑道：“我们不过当丫头一 辈子罢咧，有什么造化呢。”}
}
{
    \eR{[English source not present in the checked-in Bencraft/Joly files.]}
}
{
}

\Verse{298}
{
    \zR{宝玉笑道：“果然能够一辈子是丫头，你这个造化比我 们还大呢。”}
}
{
    \eR{[English source not present in the checked-in Bencraft/Joly files.]}
}
{
}

\Verse{299}
{
    \zR{莺儿听见这话，似乎又是疯话了，恐怕自己招出宝玉的病根来，打算 着要走。}
}
{
    \eR{[English source not present in the checked-in Bencraft/Joly files.]}
}
{
}

\Verse{300}
{
    \zR{只见宝玉笑着说道：“傻丫头，我告诉你罢。”}
}
{
    \eR{[English source not present in the checked-in Bencraft/Joly files.]}
}
{
}

\Verse{301}
{
    \zR{未知宝玉又说出什么话来，且听下回分解。}
}
{
    \eR{[English source not present in the checked-in Bencraft/Joly files.]}
}
{
}

\Verse{302}
{
    \zR{(本章完)}
}
{
    \eR{[English source not present in the checked-in Bencraft/Joly files.]}
}
{
}
